%%%%%%%%%%%%%%%%%%%%%%%%%%%%%%%%%%%%%%%%%%%%%%%%%%%%%%%%%%%%%%%
%% Section 3 -- Motivating Experiment
%%%%%%%%%%%%%%%%%%%%%%%%%%%%%%%%%%%%%%%%%%%%%%%%%%%%%%%%%%%%%%%

\section{Motivating Experiment}
\label{sec:motivating}

Before committing to the full \textbf{ACR-VLA} pipeline, we ask a targeted question: \emph{does per-DOF uncertainty carry more actionable information for correction than a scalar summary?}  We design a controlled 3D point-reaching experiment that isolates this factor.

\paragraph{Setup.}
A synthetic policy outputs actions $a\in\mathbb{R}^{3}$ with known anisotropic Gaussian noise: $\sigma_x{=}0.5,\;\sigma_y{=}0.1,\;\sigma_z{=}0.1$, simulating the depth ambiguity a monocular camera introduces along the $x$-axis.  At test time a target shift of $+0.3$ is applied along $x$, creating a systematic perturbation whose correction should be modulated by uncertainty structure.

We train three correction MLPs (2 layers, hidden dim 64, 20{,}000 training samples), differing only in their uncertainty conditioning:
\begin{itemize}[nosep]
  \item \textbf{Model A} (scalar): conditioned on the average standard deviation $\bar{\sigma}=\tfrac{1}{3}\sum_j\sigma_j$.
  \item \textbf{Model B} (per-DOF): conditioned on the full vector $[\sigma_x,\sigma_y,\sigma_z]$.
  \item \textbf{Model C} (none): conditioned on observation features only, with no uncertainty input.
\end{itemize}
To match the regime of our full pipeline, where ground-truth variances are unavailable, we report results with \emph{estimated} uncertainty computed from $K{=}10$ samples (the same budget used in \textbf{ACR-VLA}).

\paragraph{Results.}
Table~\ref{tab:motivating} reports per-axis MSE over 5{,}000 test episodes at a 5:1 anisotropy ratio ($\sigma_x/\sigma_y$) with $K{=}10$ estimated uncertainty.

\begin{table}[t]
\centering
\caption{Per-axis MSE ($\downarrow$) of correction models on the 3D point-reaching task (5:1 anisotropy, $K{=}10$ estimated uncertainty). Model~B, which receives per-DOF uncertainty, preserves accuracy on the low-uncertainty axes ($y,z$) where Model~A and~C over-correct.}
\label{tab:motivating}
\small
\begin{tabular}{lccccc}
\toprule
Model & Conditioning & MSE$_x$ & MSE$_y$ & MSE$_z$ & Total MSE \\
\midrule
A & scalar $\bar\sigma$ & 0.1874 & 0.0247 & 0.0243 & 0.0788 \\
\textbf{B} & \textbf{per-DOF} $[\sigma_x,\sigma_y,\sigma_z]$ & \textbf{0.1829} & \textbf{0.0191} & \textbf{0.0185} & \textbf{0.0735} \\
C & none & 0.2151 & 0.0250 & 0.0246 & 0.0882 \\
\bottomrule
\end{tabular}
\end{table}

Model~B reduces total MSE by 6.7\% relative to~A and by 16.7\% relative to~C.  The gains concentrate on the non-perturbed axes: $y$- and $z$-axis MSE drop by 22.7\% and 23.9\% versus~A, indicating that Model~B learns to leave low-uncertainty dimensions largely untouched while directing correction effort toward the high-uncertainty $x$-axis.  Model~A, receiving only a scalar average, lacks the information to make this distinction and applies a uniform correction that degrades the already-accurate axes.

With oracle (ground-truth) variances, the ranking $\text{B}>\text{A}>\text{C}$ is preserved with larger margins (non-perturbed axis improvements of 30--38\% over~A).  Moving from oracle to $K{=}10$ estimation reduces relative gains by roughly 15 percentage points, reflecting finite-sample noise, but does not alter the ordering.  The pattern also holds at weaker anisotropy ratios (2:1 and 1.5:1): across all six conditions tested, Model~B consistently outperforms both baselines.

\paragraph{Takeaway.}
Even a simple MLP can exploit per-DOF uncertainty structure for dimension-selective correction---a capability that scalar summaries cannot provide.  This motivates the design of \textbf{ACR-VLA}, which extracts rich per-DOF features from a flow-matching VLA and feeds them into a learned correction module (Section~\ref{sec:method}).


%%%%%%%%%%%%%%%%%%%%%%%%%%%%%%%%%%%%%%%%%%%%%%%%%%%%%%%%%%%%%%%
%% Section 4 -- Method
%%%%%%%%%%%%%%%%%%%%%%%%%%%%%%%%%%%%%%%%%%%%%%%%%%%%%%%%%%%%%%%

\section{Method}
\label{sec:method}

\textbf{ACR-VLA} (Adaptive Conformal Refinement for VLA) is an inference-time, post-hoc module that attaches to any flow-matching VLA without modifying its weights.  It comprises three components (Figure~\ref{fig:pipeline}):
\begin{enumerate}[nosep]
  \item \textbf{Conformalized Gaussian Scoring} (CGS, \S\ref{sec:cgs}): draws $K{=}10$ action samples per step, fits a shrinkage-regularized Gaussian, and extracts a 40-dimensional structured uncertainty feature.
  \item \textbf{Adaptive Conformal Inference with Chunk Overlap Consistency} (ACI+COC, \S\ref{sec:aci}): maintains an online conformal threshold that adapts to distribution shift using a self-referential coverage signal from overlapping action chunks.
  \item \textbf{Conformal Refinement Module} (CRM, \S\ref{sec:crm}): a lightweight MLP that takes the CGS features as conditioning and outputs a dimension-selective residual correction.
\end{enumerate}

We instantiate \textbf{ACR-VLA} on SmolVLA~\cite{shukor2025smolvla}, a 450M-parameter flow-matching VLA that predicts action chunks of size 50 with action dimension $d{=}6$ (3 translation + 3 rotation; the gripper is handled separately outside the conformal pipeline).


%%% 4.1 CGS %%%
\subsection{Conformalized Gaussian Scoring}
\label{sec:cgs}

CGS converts a set of stochastic action samples from the base VLA into a compact uncertainty descriptor that captures \emph{which} degrees of freedom are uncertain and \emph{how} they co-vary---information that a scalar nonconformity score discards.

\paragraph{Sampling and truncation.}
At each control step, we draw $K{=}10$ action chunks from the frozen base VLA via independent flow-matching rollouts.  Although SmolVLA predicts 50-step chunks, under receding-horizon execution only the first few steps influence the robot before re-planning.  We therefore retain only the first $H{=}5$ steps, yielding $K$ samples $\{a_h^{(k)}\}_{k=1}^{K}$ with $a_h^{(k)}\in\mathbb{R}^{d}$ for each horizon step $h=1,\dots,H$.

\paragraph{Shrinkage covariance estimation.}
For each step $h$, we compute the sample mean $\boldsymbol{\mu}_h = \frac{1}{K}\sum_{k=1}^K a_h^{(k)}$ and the Ledoit--Wolf shrinkage covariance~\cite{ledoit2004well}:
\begin{equation}
  \boldsymbol{\Sigma}_h = (1-\alpha_{\mathrm{LW}})\,\hat{\boldsymbol{\Sigma}}_{\mathrm{emp}} + \alpha_{\mathrm{LW}}\,\frac{\mathrm{tr}(\hat{\boldsymbol{\Sigma}}_{\mathrm{emp}})}{d}\,\mathbf{I},
  \label{eq:shrinkage}
\end{equation}
where $\hat{\boldsymbol{\Sigma}}_{\mathrm{emp}}$ is the empirical covariance and $\alpha_{\mathrm{LW}}$ is the analytically optimal shrinkage intensity.  With a dimension-to-sample ratio $p/n = d/K = 0.6$, the eigenvalues of the raw sample covariance suffer 30--40\% estimation error due to the Marchenko--Pastur effect; shrinkage reduces this to below 5\%~\cite{ledoit2004well}.

\paragraph{Mahalanobis nonconformity score.}
We define a joint multivariate nonconformity score:
\begin{equation}
  s_h(a) = (a - \boldsymbol{\mu}_h)^\top \boldsymbol{\Sigma}_h^{-1} (a - \boldsymbol{\mu}_h).
  \label{eq:mahal}
\end{equation}
This single score captures cross-DOF covariance structure.  In contrast, treating each DOF independently and applying Bonferroni correction across all $d{\times}H = 30$ marginals yields intervals that are substantially wider than necessary.  Our joint formulation requires Bonferroni correction only across $H{=}5$ horizon steps (target level $\alpha/H$ per step), preserving tight per-step regions~\cite{braun2025cgs}.

\paragraph{CGS feature vector.}
We assemble a 40-dimensional feature $\boldsymbol{\varphi}_{\mathrm{CGS}}$ that is passed to the refinement module:
\begin{equation}
  \boldsymbol{\varphi}_{\mathrm{CGS}} = \bigl[\,\underbrace{\sigma_{j,h}}_{j=1..d,\;h=1..H}\;;\;\underbrace{q_h}_{h=1..H}\;;\;\underbrace{v_h}_{h=1..H}\,\bigr] \;\in\;\mathbb{R}^{40},
  \label{eq:cgs_feat}
\end{equation}
where $\sigma_{j,h} = \sqrt{(\boldsymbol{\Sigma}_h)_{jj}}$ is the per-DOF standard deviation (30D), $q_h$ is the ACI-calibrated Mahalanobis threshold at step $h$ (5D, see \S\ref{sec:aci}), and $v_h = \tfrac{1}{2}\log\det(\boldsymbol{\Sigma}_h)$ is the log-volume of the covariance ellipsoid (5D).  Together, these encode \emph{which dimensions are uncertain} (per-DOF std), the \emph{overall calibrated confidence level} (conformal radius), and the \emph{spread of the predictive distribution} (log-volume).  Compared with the scalar Euclidean distance used in ReconVLA~\cite{chen2026reconvla}, this representation captures substantially more structure.

\paragraph{Relationship to single-pass uncertainty.}
Recent work~\cite{xing2026diu} shows that for flow-matching models the divergence of the velocity field recovers per-DOF aleatoric uncertainty from a single forward pass.  CGS instead estimates \emph{model uncertainty} via multi-sample covariance: a sampling-based signal that reflects both epistemic and aleatoric components, and that rises under distribution shift when the policy is uncertain about the correct action.  For action correction, this shift-sensitive signal is more actionable than a purely aleatoric estimate, since it identifies \emph{where the model needs help} rather than where the task is inherently ambiguous.  We validate this distinction empirically in ablation~A7.


%%% 4.2 ACI + COC %%%
\subsection{Adaptive Conformal Inference via Chunk Overlap Consistency}
\label{sec:aci}

CGS requires a conformal threshold $q_h$ to convert Mahalanobis scores into calibrated prediction regions.  Static calibration on a held-out set, as in conformal quantile regression~\cite{romano2019cqr}, assumes exchangeability between calibration and deployment data---an assumption that breaks under the distribution shifts common in real-world manipulation.  We instead adopt Adaptive Conformal Inference (ACI)~\cite{gibbs2021aci}, which provides long-run coverage guarantees under arbitrary distribution shift, and equip it with a novel self-referential coverage signal.

\paragraph{The coverage signal problem.}
ACI updates the miscoverage level $\alpha_t$ at each step based on whether the true action falls inside the conformal region.  In robotics, the most direct signal would check whether the \emph{executed} action lies within the predicted region.  However, contact dynamics cause the achieved end-effector displacement to differ from the commanded action, corrupting this signal.

\paragraph{Chunk Overlap Consistency (COC).}
We exploit the temporal overlap inherent in receding-horizon action chunking.  At control step $t$, the VLA predicts a chunk $[a_t^1,\dots,a_t^H]$ with associated conformal ellipsoids $\{C_t^h\}_{h=1}^H$ derived from the Mahalanobis score (\ref{eq:mahal}) and threshold $q_h$.  After executing $a_t^1$ and observing the new state $s_{t+1}$, we draw $K{=}10$ fresh samples and compute the new first-step mean $\boldsymbol{\mu}_{t+1}^1$.  The COC coverage indicator is:
\begin{equation}
  \mathrm{err}_t = \mathbb{1}\bigl[\boldsymbol{\mu}_{t+1}^1 \notin C_t^2\bigr],
  \label{eq:coc}
\end{equation}
i.e., we check whether the next step's predicted action falls within the \emph{previous} step's conformal region for the overlapping horizon position.

COC has four properties that make it well-suited for online calibration:
(i)~it is self-referential, requiring no external supervisor or trained reward model;
(ii)~it compares prediction-to-prediction rather than prediction-to-execution, making it robust to contact dynamics;
(iii)~it directly measures calibration quality---whether the conformal regions are appropriately sized;
(iv)~it produces a signal at every control step, not only at episode boundaries.

\paragraph{Online update rule.}
The ACI update follows \cite{gibbs2021aci}:
\begin{equation}
  \alpha_{t+1} = \alpha_t + \gamma\bigl(\alpha_{\mathrm{target}} - \mathrm{err}_t\bigr),
  \label{eq:aci_update}
\end{equation}
where $\gamma$ is the learning rate and $\alpha_{\mathrm{target}}$ is the desired miscoverage rate.  When $\mathrm{err}_t{=}1$ (the region failed to cover), $\alpha_{t+1}$ decreases, widening the next region; when $\mathrm{err}_t{=}0$, $\alpha_{t+1}$ increases, tightening it.  Within each episode, we apply exponentially weighted conformal prediction~\cite{barber2023conformal} with decay parameter $\rho$ to down-weight stale calibration data.  Across episodes, the fraction of intra-episode COC hits serves as $\mathrm{err}_t$ for a cross-episode ACI update, enabling the threshold to adapt over the course of deployment.


%%% 4.3 CRM %%%
\subsection{Conformal Refinement Module}
\label{sec:crm}

The Conformal Refinement Module (CRM) is a lightweight network that performs dimension-selective action correction, conditioned on the structured uncertainty features from CGS.

\paragraph{Architecture.}
CRM is a 2-layer MLP with hidden dimension 256 (${\sim}$337K parameters).  Its input concatenates three sources:
(i)~the frozen observation embedding $\mathbf{z}_{\mathrm{obs}}$ from the base VLA encoder,
(ii)~the 40D CGS feature $\boldsymbol{\varphi}_{\mathrm{CGS}}$, and
(iii)~the mean action $\boldsymbol{\mu} \in \mathbb{R}^{30}$ ($d{\times}H = 6{\times}5$).
The output is a residual correction $\boldsymbol{\delta} \in \mathbb{R}^{30}$, and the refined action is:
\begin{equation}
  \mathbf{a}^* = \boldsymbol{\mu} + \mathrm{CRM}(\mathbf{z}_{\mathrm{obs}},\;\boldsymbol{\varphi}_{\mathrm{CGS}},\;\boldsymbol{\mu}).
  \label{eq:crm}
\end{equation}

\paragraph{Training.}
CRM is trained with a single objective:
\begin{equation}
  \mathcal{L} = \bigl\|\mathbf{a}^* - \mathbf{a}_{\mathrm{expert}}\bigr\|_2^2,
  \label{eq:crm_loss}
\end{equation}
on demonstration trajectories from LIBERO~\cite{liu2023libero} augmented with synthetic perturbations (camera shifts, sensor noise, robot-state perturbation, lighting and background variation).  During training, CGS features are computed online: for each augmented observation, we run $K{=}10$ forward passes through the frozen base VLA and extract $\boldsymbol{\varphi}_{\mathrm{CGS}}$.  This ensures CRM sees realistic uncertainty patterns that co-vary with the perturbation type---for instance, camera shifts produce elevated translational variance while leaving rotational variance largely unchanged.

ACI operates exclusively at inference time: it adjusts the conformal threshold $q_h$, which enters CRM through the conformal-radius components of $\boldsymbol{\varphi}_{\mathrm{CGS}}$.  This decoupling of training and calibration eliminates the multi-objective conflicts that arise when calibration-aware losses compete with the primary correction objective in a small-capacity network~\cite{wang2025cploop}.

\paragraph{Why per-DOF conditioning enables selective correction.}
To build intuition, consider a camera perturbation that introduces translational ambiguity in the $x$-$y$ plane while leaving $z$-translation and rotation well-determined.  A scalar nonconformity score (e.g., average Euclidean sample distance, as in \cite{chen2026reconvla}) compresses this into ``moderate uncertainty everywhere,'' providing no basis for selective correction.  The per-DOF structure in $\boldsymbol{\varphi}_{\mathrm{CGS}}$ instead reveals that $\sigma_{x},\sigma_{y} \gg \sigma_{z},\sigma_{\mathrm{rot}}$, enabling CRM to apply large corrections to the uncertain translational axes while preserving the base VLA's accurate rotational predictions.  Section~\ref{sec:motivating} validates this mechanism in a controlled setting; Section~\ref{sec:experiments} demonstrates it on real VLA tasks.

\paragraph{Inference protocol.}
Each control step proceeds as: draw $K{=}10$ samples $\to$ compute CGS $\to$ apply CRM $\to$ execute the first action of the refined chunk $\to$ observe new state $\to$ draw $K{=}10$ fresh samples $\to$ compute COC (\ref{eq:coc}) $\to$ ACI update (\ref{eq:aci_update}).  The entire module adds ${\sim}$337K parameters and $K{=}10$ extra forward passes to the base VLA, with no retraining of the base model.
